%%%%%%%%%%%%%%%%%%%%%%%%%%%%%%%%%%%%%%%%%%%%%%%%%%%%%%%%%%%%%%%%%%%%%%%%
% Plantilla TFG/TFM
% Escuela Politécnica Superior de la Universidad de Alicante
% Realizado por: Jose Manuel Requena Plens
% Contacto: info@jmrplens.com / Telegram:@jmrplens
%%%%%%%%%%%%%%%%%%%%%%%%%%%%%%%%%%%%%%%%%%%%%%%%%%%%%%%%%%%%%%%%%%%%%%%%

\chapter{Introducción}

\section{Importancia del PIB}

El Producto Interior Bruto (PIB) es uno de los principales indicadores macroeconómicos utilizados para medir el nivel de actividad económica de un país. Este indicador recoge el valor monetario total de los bienes y servicios finales producidos dentro de una economía durante un periodo de tiempo determinado, generalmente un trimestre o un año. Debido a su carácter agregado, el PIB se ha consolidado como una referencia fundamental para analizar el desempeño económico de los países y evaluar la evolución de su actividad productiva.

Su evolución proporciona información clave sobre el crecimiento económico, los ciclos económicos y la salud general de una economía. A través del análisis de las variaciones del PIB es posible identificar periodos de expansión, desaceleración o recesión económica, así como evaluar el impacto de distintos factores económicos y políticos sobre el conjunto del sistema productivo.

La capacidad de anticipar cambios en el PIB resulta especialmente relevante para gobiernos, bancos centrales y agentes económicos, ya que permite una mejor planificación de políticas económicas y una toma de decisiones más informada. Por ejemplo, las autoridades monetarias utilizan las previsiones de crecimiento económico como una de las variables fundamentales para determinar la orientación de la política monetaria, mientras que los gobiernos emplean estas estimaciones para planificar políticas fiscales y presupuestarias.

Además de su uso como indicador del crecimiento económico, el PIB es una variable de referencia fundamental para instituciones nacionales e internacionales. Organismos como bancos centrales, gobiernos y entidades supranacionales utilizan la evolución del PIB para evaluar el ciclo económico, diseñar políticas económicas y comparar el desempeño de distintas economías.

En contextos como la Unión Europea, el análisis comparativo del PIB entre países resulta especialmente relevante, ya que permite identificar diferencias estructurales, procesos de convergencia económica y efectos derivados de políticas económicas comunes. Este tipo de análisis resulta esencial para comprender la evolución económica relativa de los distintos Estados miembros y para evaluar la efectividad de determinadas medidas económicas a nivel supranacional.

Por ello, disponer de herramientas que permitan analizar y anticipar la evolución del PIB adquiere una importancia estratégica tanto a nivel nacional como internacional. Una comprensión adecuada de las cuentas nacionales y de los indicadores asociados a la actividad económica resulta fundamental para el análisis macroeconómico moderno y para la correcta interpretación de la dinámica económica de los países. \citet{oecd_understanding_national_accounts_2014}

\section{Dificultad de la predicción económica}

La predicción de variables macroeconómicas presenta importantes desafíos debido a la naturaleza compleja y cambiante de los sistemas económicos. A diferencia de otros fenómenos en los que las relaciones entre variables pueden modelizarse mediante leyes relativamente estables, las economías están influenciadas por una gran variedad de factores institucionales, sociales y políticos que pueden modificar significativamente su comportamiento a lo largo del tiempo.

Factores internos y externos, como crisis financieras, cambios en la política económica o perturbaciones globales, pueden alterar de forma significativa la evolución del PIB. Eventos inesperados, como crisis económicas internacionales, conflictos geopolíticos o pandemias, pueden provocar cambios bruscos en la actividad económica que resultan difíciles de anticipar mediante modelos tradicionales.

A esta complejidad se suma el carácter no lineal de muchas relaciones económicas, así como la presencia de incertidumbre y eventos inesperados que dificultan la modelización precisa del comportamiento económico. Además, las economías presentan estructuras heterogéneas, con distintos niveles de desarrollo, especialización productiva y marcos institucionales, lo que complica la aplicación de modelos universales válidos para distintos países o periodos temporales.

Otro elemento relevante es la limitada disponibilidad de datos en tiempo real y las posibles revisiones posteriores de las estadísticas macroeconómicas. Las estimaciones iniciales del PIB y de otros indicadores económicos suelen revisarse a medida que se dispone de información más completa, lo que introduce un grado adicional de incertidumbre en los procesos de predicción económica.

En este contexto, la predicción del crecimiento económico se convierte en un problema especialmente complejo que combina incertidumbre estructural, limitaciones en la información disponible y posibles cambios en las relaciones entre variables económicas a lo largo del tiempo. \citet{stock_watson_2003}

Estas dificultades ponen de manifiesto la necesidad de explorar enfoques alternativos que permitan capturar patrones complejos en los datos y mejorar la capacidad predictiva frente a los métodos tradicionales.

\section{IA y reconocimiento de patrones en datos económicos}

En los últimos años, los avances en inteligencia artificial y aprendizaje automático han permitido desarrollar modelos capaces de identificar patrones complejos en grandes conjuntos de datos. Estas técnicas han demostrado su utilidad en numerosos ámbitos, como la visión artificial, el procesamiento del lenguaje natural o el análisis de grandes volúmenes de información, y su aplicación al análisis económico ha despertado un creciente interés tanto en el ámbito académico como en el profesional.

A diferencia de los enfoques econométricos clásicos, muchas técnicas de aprendizaje automático no requieren supuestos estrictos sobre la forma funcional de las relaciones entre variables. En lugar de definir previamente una estructura matemática concreta, estos modelos aprenden dichas relaciones directamente a partir de los datos disponibles, lo que les permite adaptarse a estructuras más complejas y dinámicas.

Su capacidad para analizar grandes volúmenes de datos y detectar patrones subyacentes las convierte en herramientas especialmente adecuadas para el estudio de fenómenos económicos, donde las relaciones entre variables pueden ser no lineales y estar sujetas a cambios a lo largo del tiempo. \citet{mullainathan_spiess_2017}

En este sentido, el uso de modelos basados en inteligencia artificial puede complementar los métodos tradicionales de análisis económico. Mientras que los enfoques econométricos clásicos se centran principalmente en la interpretación causal de las relaciones entre variables, muchas técnicas de aprendizaje automático priorizan la capacidad predictiva y la identificación de patrones en los datos.

Esta complementariedad ha motivado un creciente número de investigaciones que exploran la aplicación de técnicas de aprendizaje automático al análisis macroeconómico, incluyendo la predicción de variables agregadas como el crecimiento del PIB. \citet{pascual_ruiz_ml_series_2021}

\section{Objetivos}

A nivel personal, como estudiante del doble grado en Ingeniería Informática y Administración y Dirección de Empresas, este trabajo surge del interés por la aplicación de técnicas y modelos de inteligencia artificial al análisis económico. En particular, se pretende explorar cómo las economías responden ante distintos cambios a través del estudio de datos macroeconómicos, profundizando en el uso de técnicas de aprendizaje automático y en su capacidad para identificar patrones relevantes en este contexto.

El objetivo principal de este Trabajo de Fin de Grado es analizar y predecir la evolución del Producto Interior Bruto mediante técnicas de análisis de datos y aprendizaje automático, utilizando información macroeconómica de distintos países y evaluando el comportamiento de los modelos en diferentes escenarios.

Para ello, se plantea el desarrollo de un proceso completo que incluye la recopilación y preparación de datos macroeconómicos, la construcción de un conjunto de datos estructurado, el entrenamiento de modelos predictivos y la evaluación de su capacidad de generalización. Asimismo, se pretende analizar cómo la inclusión de distintas variables macroeconómicas puede influir en el rendimiento de los modelos de predicción.

Finalmente, como complemento al análisis realizado, se desarrolla una aplicación web interactiva que permite visualizar los resultados de los modelos, explorar las predicciones obtenidas y facilitar la interpretación de los experimentos realizados. Este componente busca aportar una dimensión más aplicada al trabajo, permitiendo representar de forma clara los resultados obtenidos durante el proceso de modelización.